\begin{lemma}
Let \(\varepsilon>0\).  Then, for all sufficiently large \(n\), with
\(L=\log n\), both strengthened absorption estimates
\[
36e(1+\varepsilon)L^3\le n^{2\varepsilon}
\]
and
\[
24L^2\le (1+\varepsilon)\sqrt L\,n^\varepsilon
\]
hold.
\end{lemma}
\begin{proof}
Since \(\varepsilon>0\), also \(2\varepsilon>0\) and \(1+\varepsilon>0\).
The standard logarithm-versus-power estimate, formalized in Mathlib as
\(\log x=o(x^a)\) for every \(a>0\), implies that every fixed positive power
of \(\log n\) is eventually dominated by every fixed positive power of \(n\).
Applying this with the positive exponent \(2\varepsilon\) and the positive
constant \(1/(36e(1+\varepsilon))\), there is \(N_1\) such that for all
\(n>N_1\),
\[
(\log n)^3\le \frac{1}{36e(1+\varepsilon)}\,n^{2\varepsilon}.
\]
Multiplying by the positive constant \(36e(1+\varepsilon)\) gives
\[
36e(1+\varepsilon)(\log n)^3\le n^{2\varepsilon}. \tag{1}
\]

Apply the same logarithm-versus-power estimate with the positive exponent
\(\varepsilon\) and the positive constant \((1+\varepsilon)/24\).  There is
\(N_2\) such that for all \(n>N_2\),
\[
(\log n)^2\le \frac{1+\varepsilon}{24}\,n^\varepsilon,
\]
and therefore
\[
24(\log n)^2\le (1+\varepsilon)n^\varepsilon. \tag{2}
\]
Finally enlarge the threshold to \(N_3\) so that \(\log n\ge1\) whenever
\(n>N_3\).  Then \(\sqrt{\log n}\ge1\), so multiplying the right side of
(2) by \(\sqrt{\log n}\) can only increase it:
\[
(1+\varepsilon)n^\varepsilon
\le (1+\varepsilon)\sqrt{\log n}\,n^\varepsilon.
\]
Combining this inequality with (2) gives
\[
24(\log n)^2\le (1+\varepsilon)\sqrt{\log n}\,n^\varepsilon. \tag{3}
\]
Taking \(N=\max(N_1,N_2,N_3)\), the two desired estimates are exactly (1)
and (3) for every \(n>N\).
\end{proof}
